\section{Experimental Results}
\label{Exp}
\subsection{Experiment Setup}
\label{Exp:Set}
We compare our methodology ({\it Proposed}) against two methods: Qiskit's \texttt{PhaseOracle} method
({\it PhaseOracle}), and the method proposed in Meuli et al~\cite{bib-meuli-esop} implemented in the
Caterpillar library~\cite{bib-meuli-ros} ({\it Caterpillar}).

The {\it Proposed} method implements the method in~\secref{Pro}. It takes an ESOP expression and
feeds it into EXORCISM-4 as an optimizer. EXORCISM-4 was chosen for its lightweight implementation.
The logic network mapping is accomplished using abc~\cite{bib-mishchenko-lut} for 2-fanin nodes,
and using the \texttt{mfs} feature in order to optimize for number of nodes.

The \texttt{PhaseOracle} method takes in an ESOP expression and generates a set of controlled-Z gates
just as in~\secref{Pre:OracleEsop}. Any controlled-Z gates are implemented as MCT gates conjugated by
Hadamard gates as detailed in~\figref{fig-phase-n4}. The MCT gates are implemented using the technique
in~\cite{bib-rtof-maslov}, the lowest known T-count for the MCT without using measurements. Just like the
{\it Proposed} method, it optimizes the ESOP using EXORCISM-4.
Comparing {\it Proposed} to this method is meant to highlight the power of using the generated
3-bit function library over merely synthesizing using the ESOP method.

The {\it Caterpillar} method decomposes the Boolean function into an XAG graph, and optimizes
this graph to minimize for T-count and T-depth~\cite{bib-meuli-esop,bib-meuli-xag,bib-meuli-ros,bib-meuli-mult}.
To minimize the XAG graph for T-depth, we use the \texttt{xag\_minmc\_resynthesis} function in the
associated \texttt{mockturtle} library which implements the function to optimize for
T-depth in~\cite{bib-meuli-mult}. The XAG is then synthesized as a quantum circuit using the method
in~\cite{bib-meuli-mult}. Because the entire quantum oracle is a single output function, we create a phase
oracle from it by conjugating it with Hadamard gates as demonstrated in~\figref{fig-phase-n4}. This method
is included to show that directly synthesizing phase gates has significant advantages in T-count and
T-depth over using quantum oracle based synthesis.

For data, while synthesis methods for quantum oracles are often tested on reversible benchmarks, these often
represent multibit Boolean mappings. Phase oracles, on the other hand, are single output Boolean functions.
Thus, this work uses randomly generated Boolean functions to do its comparisons. For better control of
the methods' behavior, ESOP expressions are used as benchmarks to feed into all of the comparison
methods.

To provide some assurance that T-count for each of the generated circuits is reasonably optimized, all
generated circuits are additionally processed by T-par~\cite{bib-amy-matroid}.

The main code is run in Python, while T-par is executed as a system call to a C++ binary.

Here we propose two kinds of ESOP expressions to test:

\begin{itemize}
\item (N3): Randomly generated ESOP expressions are constrained to those of $n \leq 3$. This tests the
  ability of the precomputed 3-bit library to enable simplification of pure CNOT+T circuits through
  easier sum-of-paths calculations.
\item (N4): Randomly generated ESOP cubes are constrained as follows: those with up to two literals
  (i.e., no T gates) are generated with 2/5 probability, three, literals (has T-gates but CNOT+T) with
  2/5 probability, and four or more literals with 1/5 probability. This tests the ability of the proposed
  method to find a simplified solution in its EXORCISM-4 step. It creates an expression which can
  already be mostly expressed as a CNOT+T network due to its cubes of three or fewer literals, in order to
  create a more fair comparison between Qiskit's method and the proposed method. The probabilities
  are chosen so that cubes with four or more literals have as little common literals to other cubes as possible,
  reducing the possibility that the cubes of four or more literals will be reduced down by the ESOP step.
\end{itemize}

We test these for $n \leq 12$ literals. The ESOP expressions are randomly generated according to two groups based
on the number of cubes $k$ in them. These are further divided into two groups:

\begin{itemize}
\item ($k=2^{n-1}$): This randomizes the number of cubes between $2n$ and $2^{n-1}$. This increases the chances that
  inserted subcircuits will have common literals that allow them to be simplified and have reduced T-count, while
  making it harder for T-par to partition any cubes $n \geq 4$ for T-depth scheduling.
\item ($k=n$): This randomizes the number of cubes between $1$ and $2n-1$. This makes the arrangements sparser,
  allowing for better parallelization by T-par, but making it harder to have common literals that can combine to
  reduce the number of cubes.
\end{itemize}

For each combination of N3/N4 and $k=2^{n-1}$/$k=n$, we generated 100 random ESOP expressions and fed them
to all three methods. We take the T-count of the circuit generated by {\it Phase Oracle} or
{\it Caterpillar} $tcnt_c$ and the T-count of the circuit generated by Proposed $tcnt_p$ and take the
percentage difference $\Delta$T-count\% $= (tcnt_q - tcnt_p) / tcnt_q$ and take the arithmetic average of
$\Delta$T-count\% for Avg $\Delta$ T-count\%. We do the same for T-depth and get Avg $\Delta$T-depth\%.
The results are shown in Table~\ref{table-results}.
%\subsection{Results}
%\label{Exp:Res}

\begin{table}[t]
  \begin{center}
    \scalebox{1.0} {
      \begin{tabular}{c|c|c|c|c}\hline
                                   &                       &                       & $k=2^{n-1}$ & $k=n$    \\\hline
  \multirow{4}{*}{vs. PhaseOracle} &  \multirow{2}{*}{N3}  & Avg $\Delta$T-count\% &  0.328      & 0.330    \\\cline{3-5}
                                   &                       & Avg $\Delta$T-depth\% &  0.588      & 0.603    \\\cline{2-5}
                                   &  \multirow{2}{*}{N4}  & Avg $\Delta$T-count\% &  0.491      & 0.425    \\\cline{3-5}
                                   &                       & Avg $\Delta$T-depth\% &  0.633      & 0.591    \\\hline
  \multirow{4}{*}{vs. Caterpillar} &  \multirow{2}{*}{N3}  & Avg $\Delta$T-count\% &  0.781      & 0.802    \\\cline{3-5}
                                   &                       & Avg $\Delta$T-depth\% &  0.790      & 0.843    \\\cline{2-5}
                                   &  \multirow{2}{*}{N4}  & Avg $\Delta$T-count\% &  0.448      & 0.593    \\\cline{3-5}
                                   &                       & Avg $\Delta$T-depth\% &  0.351      & 0.557    \\\hline
\end{tabular}


    }
  \end{center}
  \caption{Experimental Results}
  \label{table-results}
\end{table}

\subsection{Analysis}

As can be seen from Table~\ref{table-results}, the {\it Proposed} method outperformed both of the
comparison methods, for all of the input ranges observed.

The difference was most stark with N3. Because the entire circuit can be implemented using only
CNOT+T, many of the T-gates can combine together and go away in the {\it Proposed} method. However,
with {\it Caterpillar}, gates with three inputs will need an ancilla, and so there are fewer opportunities
to combine T-gates through resynthesis. On the other hand, {\it PhaseOracle} can similarly be implemented
entirely as a CNOT+T circuit, which is reflected by a muted advantage in T-count by the {\it Proposed} method.
This is also reflected in the {\it Proposed} method's advantage over the {\it Phase Oracle} method being close
together for both $k=2^{n-1}$ and $k=n$; whether sparse or not, a CNOT+T circuit will still be easily simplified
together. However, {\it PhaseOracle} can't combine T-gates driven by negated literals, so the difference reflects
missing those opportunities.

On the other hand N4 had a more muted advantage. {\it Proposed} performs slightly better against {\it PhaseOracle}
in this set of inputs, because there are more $4+$ input gates, which benefit from sharing fanouts, whereas
{\it PhaseOracle} doesn't utilize it. For the same reason, the advantage of {\it Proposed} over {\it Caterpillar}
with the same set of inputs is lessened: with more $4+$ input gates, {\it Caterpillar} is able to match
the savings from logic sharing that {\it Proposed} is able to attain. The advantage that {\it Proposed} has over
{\it Caterpillar} in this set of inputs is from its 3-input phase oracle library, as well as needing to synthesize
a smaller function, as mentioned in~\secref{phase-direct}.

%% The most significant result is \{$k=2^{n-1}$,N4\}, which reports reductions of over 60\% compared
%% with Qiskit's method. This is because EXORCISM-4 can find more cubes to simplify when they are
%% denser and share more literals.

%% Next is \{$k=n$,N4\}. Despite the sparsity, EXORCISM-4 was still able to find cubes to simplify
%% and that contributes to the reduction. This is notable because the EXORCISM-4 produced ESOP
%% expressions had a lower average number of literals in each cube.

%% \{$k=2^{n-1}$,N3\} showed a more moderate but significant advantage. This is because EXORCISM-4
%% can find fewer advantages over the standard case since it is already in a form that allows
%% the T-par algorithm to work optimally. However, there is still a significant advantage because
%% the 3-bit library of functions allows a more efficient sum of paths calculation to combine
%% T-gates into Clifford phase gates.

%% Finally \{$k=n$,N3\} shows the least amount of advantage because of its sparsity. This is most
%% notable in T-depth, because T-par can already efficiently schedule the circuit produced by Qiskit.
%% However, there is still an advantage because of the sum of paths calculations.

%% Despite the clear advantages conferred by the Proposed method, there remain a few corner
%% cases where the Qiskit circuit outperformed the one generated by the Proposed method, such as
%% in the case for the expression $(x_3x_7) \oplus (~x_6x_4) \oplus (x_3~x_7~x_6) \oplus (~x_6x_5~x_3) \oplus (x_2x_4) \oplus (~x_8~x_3) \oplus (~x_7x_3) \oplus (x_1~x_8) \oplus (~x_1~x_8) \oplus (x_2x_4x_3) \oplus (x_3x_5) \oplus (~x_7x_6)$.
%% Evaluating this will require much more material than can be reasonably fit into this paper,
%% and we will leave it as a topic for future research.

%% \label{Exp:An}

